\documentclass[11pt]{article}
\usepackage[a4paper,margin=.8in]{geometry}
\usepackage{amsmath,amssymb,graphicx,booktabs}
\usepackage[colorlinks=true,urlcolor=blue,linkcolor=blue]{hyperref}
\graphicspath{{../output/heggie_stock2000/plots/}{../output/heggie_stock_calibrated2000/plots/}}
\title{Heggie Figs.~1--4 using unchanged ThreeFluidSim}
\author{HydroSim reproduction note}
\date{4 October 2026}
\begin{document}
\maketitle
\section{Scope and origin of the previous changes}
We rerun the gas-model comparisons of Heggie and Aarseth (1992),
\emph{MNRAS} \textbf{257}, 513,
\href{https://doi.org/10.1093/mnras/257.3.513}{paper},
using the public hydrostatic solver at upstream revision \texttt{82d09b8}.
The reproduction adds initialization, observation, run scripts and plots,
but does not change \texttt{ThreeFluidSim} or its evolution operators.
The paper's Figs.~1--4 are gas models with primordial-binary heating,
not its later $N$-body calculations and not binary-formation runs.

The earlier local experiment (\texttt{18cd6c4}) added three features:
\begin{enumerate}
\item \emph{Relative-dispersion heating}, motivated by the paper's Eq.~(3).
It replaced a one-component inverse dispersion by
$(U_s+U_b)^{-1/2}$ in BS heating and differentiated with respect to both
new energies. Specifically, for $S=U_s+U_b$ it used
$S^{-1/2}\simeq3/(2\sqrt{S^n})-S^{n+1}/[2(S^n)^{3/2}]$.
This requires cross-component matrix entries and a matching RHS; it is
not just a constant conversion factor.
\item \emph{A reflecting sphere}, specified by the paper in Sec.~3.1.
The added boundary held the outer radius fixed, set the last conduction
row to $U_{N-1}-U_{N-2}=0$, and preserved the wall's entropy-based
pressure during remapping. In the existing nonconservative conduction
stencil, this is a zero-gradient condition, not a proof of an exact
zero-flux energy ledger.
\item \emph{Configurable initial endpoints}, a numerical choice made for
the experiment. It extended the canonical Plummer initializer from its
built-in $10^{-2}$--$10^3$ domain to $10^{-6}$--$10^4$. The paper does
not prescribe these particular endpoints.
\end{enumerate}
The public hydrostatic solver has none of these options. We therefore
retain its native grid, single-component heating denominators, outer
hydrostatic equation and fixed outer $U$ during conduction. The last
radius may move. The previous local results are controls, not results
produced by this PR.

\section{What the unchanged solver implements}
All internal quantities are dimensionless; $P_i=2\rho_iU_i/3$ and
$\sigma_i^2=2U_i/3$, with $\gamma=5/3$ for both species. A binary's
$U_b$ is its centre-of-mass kinetic energy, not its binding energy.
We take $C=0.104$, $m_b=2m_s$ in the two-component runs and
$\ln\Lambda=10$. Formation, destruction, ejection and tidal removal
are off. The precise Fig.~4 Coulomb logarithm and distant boundary
location are not uniquely specified by the paper.

Define the dimensionless coefficients obtained from Eqs.~(3) and (10),
\begin{equation}
 B=\frac{3.807}{16\pi\ln\Lambda},\qquad
 D=\frac{5\sqrt3}{24\pi\ln\Lambda}.
\end{equation}
Heggie's desired specific heating terms are
\begin{equation}
 Q_s=\frac{2B\rho_b}{3\sqrt{U_s+U_b}},\qquad
 Q_b=\frac{B\rho_s}{3\sqrt{U_s+U_b}}+
       \frac{D\rho_b}{\sqrt{2U_b}}.
\end{equation}
We configure the \emph{existing} coefficients as
$c_{4,sb}=2B/(3\sqrt2)$, $c_{4,bs}=B/(3\sqrt2)$ and
$c_{4,bb}=D/\sqrt2$. Consequently, the actual stock-solver sources are
\begin{equation}
 Q_s^{\rm stock}=\frac{2B\rho_b}{3\sqrt{2U_s}},\qquad
 Q_b^{\rm stock}=\frac{B\rho_s}{3\sqrt{2U_s}}+
       \frac{D\rho_b}{\sqrt{2U_b}}.
\end{equation}
The BS heat partition remains $2:1$, and BB self-heating is exact in
this normalization. BS heating agrees at $U_s=U_b$, but its ratio to
the desired rate is
\begin{equation}
 \frac{Q_{\rm BS}^{\rm stock}}{Q_{\rm BS}^{\rm Heggie}}
 =\sqrt{\frac{U_s+U_b}{2U_s}}.
\end{equation}
At equipartition, $U_b=U_s/2$, this is $\sqrt3/2\simeq0.866$.
Thus a constant coefficient cannot make the two models equivalent at
all times. The companion Mathematica file checks this mapping and the
dimensionless conduction and exchange coefficients,
$c_{2,i}=2\sqrt{2/3}\,C m_i/m_s$ and $c_{1,ij}=1/\sqrt{3\pi}$.

For Figs.~1--2 the single effective fluid uses the solver's existing
self-heating slot \texttt{FB}; the other two slots are negligible
identical clones and all particle masses are equal. This is a storage
choice, not a physical binary fraction of unity. It reproduces the
single-fluid source $3.807f^*\rho/[8\pi\sqrt{2U}]$ without modifying
the conduction operator. Every runner calls the normal \texttt{evolve()},
and saving is performed only by a read-only observer.

\section{Units, profile fit and diagnostics}
Let $\mathcal M$ be the initial code mass and $L=3\pi/16$. We plot
\begin{equation}
 r_{\rm plot}=Lr,\qquad
 \rho_{\rm plot}=\frac{\rho}{4\pi L^3\mathcal M},\qquad
 \tau=\frac{t_{\rm code}}{0.18\ln(10)r_h(0)^{3/2}\sqrt{\mathcal M}}.
\end{equation}
The base-ten logarithm in the paper's relaxation time, Eq.~(15), is
converted explicitly. The baseline continuum Plummer normalization has
$G=M(0)=1$, $E(0)=-1/4$.

The optional Fig.~4 fit varies the singles' Plummer scale $a$ and the
relative binary scale $b$, keeping $\mu=M_b/M_s=2f/(1-f)$ fixed at
$f=N_b/(N_s+N_b)=0.03$:
\begin{equation}
 \rho_s=\frac{(1+r^2/a^2)^{-5/2}}{(1+\mu)a^3},\qquad
 \rho_b=\frac{\mu(1+r^2/(ab)^2)^{-5/2}}{(1+\mu)(ab)^3}.
\end{equation}
Canonical initialization reconstructs pressure in the common potential.
The homology $R\mapsto aR$, $\rho\mapsto\rho/a^3$,
$P\mapsto P/a^4$, $U\mapsto U/a$ preserves discrete hydrostatic balance.
The native grid scales with $a$; the observer length unit does not change.
Neither equal component dispersions nor $E(0)=-1/4$ is imposed after
fitting different component scales.

At 2000 shells the fresh stock-solver fit gives
$a=0.7172313$, $b=1.5692616$, and
$(\rho_{c1},\rho_{c2})=(3.599988,0.600000)$ at $\tau=1$.
The targets $(3.6,0.6)$ are approximate readings of the solid paper
curves, uncertain at roughly 5--10\%; the quoted simulation digits do
not imply precise reference data. Both heating cases use this same fit.
Only that epoch is fitted; Table~\ref{tab:fit} compares later predictions.
\begin{table}[ht]
 \centering
 \input{../output/heggie_stock_calibrated2000/plots/density_table.tex}
 \caption{BS+BB central densities. Only the $\tau=1$ row constrains the fit.}
 \label{tab:fit}
\end{table}

The baseline BS+BB calculation reaches $\tau=150.040$, with singles'
peak density $5.499$ at $\tau=10.566$, versus $5.134$ in the earlier
modified-solver experiment. The large baseline Fig.~4 offset therefore
persists. Its central BS rate ratio ranges from $1$ to $0.869$.
The last radius changes by about $7.5\%$; at most $0.073\%$ of the mass
lies outside code radius 100. These are diagnostics, not proof that
the boundary is irrelevant. Its total-mass drift is $0.542\%$; the
stock remapping is not exactly conservative. Mechanical energy is not
a closed-system conservation residual because heating draws from an
unresolved binary binding-energy reservoir.

The fitted BS+BB run reaches $\tau=150.216$. Its singles peak is
$14.774$ at $\tau=10.782$, and its initial plotted energy is $-0.33946$.
At $\tau=10$, the fitted densities $14.618$ and $15.752$ exceed the
paper readings by about $35\%$ and $66\%$. The fitted BS-only run
still has no bounce and reaches the density ceiling at $\tau=12.0356$.
Thus fitting the left edge does not explain all later differences.
The fitted BS+BB mass drift is $0.708\%$. Repeating it at 1000 shells,
with the same parameters and no refit, changes central densities by at
most $1.53\%$ and $1.31\%$ on 1000 common logarithmic times from
$\tau=1$ to 150. This check does not establish exact conservation.

\section{The curves in Figs.~1--4}
\paragraph{Figure 1.}
Five colours label $f^*=0.001$, $0.002$, $0.005$, $0.01$, $0.02$,
the effective heating strength
$f^*=f(1-f)/[\ln\Lambda(1+f)^2]$ in the paper's Eq.~(23).
The vertical variable is $M_c/M(0)$, with
$M_c=0.517(4\pi/3)\rho_cr_c^3$ in physical units.
Weak heating allows continued collapse; sufficiently strong heating
produces a minimum and approximately steady core mass fraction.

\paragraph{Figure 2.}
The same colours represent the same five heating strengths. Our solid
curves are the core radius $r_c$, and dashed curves the half-mass radius
$r_h$, against $\tau$. We use
$r_c^2=9\sigma_c^2/(4\pi G\rho_c)$; $r_h$ encloses half the current
total mass. The paper uses different styles for different $f^*$,
with each labelled lower $r_c$ curve paired to an upper $r_h$ curve.

\paragraph{Figure 3.}
No heating is present. The initial density ratio is $0.01$, equal
Plummer scales give equal dispersions, and $m_b=2m_s$.
Purple is $\rho_{c2}/\rho_{c1}$; red and blue are the singles' and
binaries' growth rates
$\xi_i=t_{rc,i}\,d\ln\rho_{ci}/dt$, with
$t_{rc,i}=0.34\sigma_{ci}^3/(G^2m_i\rho_{ci}\ln\Lambda)$.
The horizontal variable is $\Phi_c/\Phi_c(0)$, not time.
The growing density ratio shows segregation of the heavy component.

\paragraph{Figure 4.}
Blue is singles' central \emph{mass} density, red binaries'.
Solid lines include BS+BB heating; dashed lines include BS only.
The same solid/dashed distinction holds in the gray paper panel.
A density-ceiling termination must not be described as a reproduced
bounce. In the separate before/after PDF all coloured curves are
BS+BB, with dashed meaning original profile and solid meaning fitted
profile instead; its legend states this different encoding.

\clearpage
\begin{figure}[p]
 \centering
 \includegraphics[width=.73\textwidth]{heggie_fig1_overlay.pdf}\\[-1ex]
 \includegraphics[width=.73\textwidth]{heggie_fig2_overlay.pdf}
 \caption{Original-profile Figs.~1--2: gray paper and coloured
 unchanged-solver histories. The five heating strengths share colours
 across the two panels. No initial-profile calibration is applied.}
\end{figure}
\begin{figure}[p]
 \centering
 \includegraphics[width=.73\textwidth]{heggie_fig3_overlay.pdf}\\[-1ex]
 \includegraphics[width=.73\textwidth]{../output/heggie_stock_calibrated2000/plots/heggie_fig4_overlay.pdf}
 \caption{Fig.~3 uses the original profile; Fig.~4 uses the freshly
 calibrated profile with unchanged ThreeFluidSim. Reference aspect ratios
 and axis limits are preserved. Gray is the paper, colour the simulation.}
\end{figure}
\clearpage
\section{Reconstruction and verification}
See \texttt{docs/heggie-stock-hydro.md} for build, CLI, data-column and
stop-condition details. Generate the baseline and fitted suites with
\texttt{run\_heggie\_hydro\_suite.py} and \texttt{calibrate\_heggie\_hydro.py},
then run \texttt{plot\_heggie\_hydro.py} and
\texttt{plot\_heggie\_calibrated.py} with a supplied copy of the paper.
Compile this source from \texttt{docs/}, directing generated TeX outputs
to \texttt{../output/}. Generated histories and PDFs are not committed.

Strict O3/native builds omit fast-math. All simulations use one BLAS and
OpenMP thread. Tests cover initialized hydrostatic balance, unchanged
native boundaries and actual heating stencil, profile mass fractions,
observer readback, refusal to overwrite runs, and plotting interpolation
without extrapolation. The full existing test suite and the Heggie checks
pass. Mathematica verifies twelve coefficient/source identities and four
profile/homology identities. No solver implementation is duplicated inside
the reproduction and no experimental solver flag is required.
\end{document}
